\section{Molecular dynamics evolution}
\label{sec:force}

The analyticity of the stout-link scheme permits the use of modern Monte
Carlo updating methods for gauge-field actions constructed using stout links.
Molecular dynamics evolution forms the core of many of the Markov processes
used to generate ensembles of field configurations needed for unquenched
QCD simulations\cite{hmcgauge,hmcfermion,hmc,Ralgorithm}.
Since a small change to the underlying gauge
fields always leads to a small change in the stout links, then provided
the small change to the stout links causes small changes to the action of
the theory, the force term on the underlying links is always well-defined. 
The importance of defining a link smearing scheme which permits the
computation of the molecular dynamics force term has been emphasized
recently in Ref.~\cite{flic2}.

The $\Sigma_\mu(x)$ force field describes how the action $S$ changes
when the gauge-field link variables $U_\mu(x)$ change, holding the momenta 
$P_\mu(x)$ and the pseudofermion field $\phi(x)$ fixed.
For this reason, we shall write the action $S=S[U]$ in this section,
ignoring its dependence on the pseudofermion field.  The transpose
of the force field is defined by the derivative of the action with
respect to the link variables:
\begin{equation}
 \Sigma_\mu(x)=\left(\frac{\partial S[U]}{\partial U_\mu(x)}\right)^T.
\end{equation}
Assume that the action can be written as a sum of
a term $S_{\rm th}[U]$ constructed entirely of the original thin
link variables and a term $S_{\rm st}[\tilde{U}]$ constructed
completely out of stout links:
\begin{equation}
  S[U]=S_{\rm th}[U]+S_{\rm st}[\tilde{U}].
\end{equation}
Then the force field may be written
\begin{equation}
  \Sigma_\mu(x)=\Sigma_\mu^{\rm(th)}(x)+\Sigma_\mu^{(0)}(x),
\end{equation}
where
\begin{equation}
 \Sigma_\mu^{\rm(th)}(x)=\left(\frac{\partial S_{\rm th}[U]}{
  \partial U_\mu(x)}\right)^T\!\!\!, \quad
 \Sigma_\mu^{\rm(0)}(x)=\left(\frac{\partial S_{\rm st}[\tilde{U}]}{
  \partial U_\mu(x)}\right)^T\!\!\!.
\end{equation}
The computation of $\Sigma_\mu^{\rm(th)}(x)$ is usually
straightforward, and nothing more about its determination needs
discussion here.  We now focus on the evaluation of 
$\Sigma_\mu^{\rm(0)}(x)$. 

Recall that the stout links are constructed iteratively starting
with the original links 
$U\rightarrow U^{(1)}\rightarrow U^{(2)}
 \rightarrow\cdots\rightarrow U^{(n_\rho)}=\tilde{U}$.  
The computation of the force field proceeds similarly in an
iterative fashion, except that the order is reversed
$\tilde\Sigma=\Sigma^{(n_\rho)}\rightarrow\Sigma^{(n_\rho\!-\!1)}
\rightarrow\cdots\rightarrow\Sigma^{(1)}\rightarrow\Sigma^{(0)}$.
The sequence starts by computing
\begin{equation}
 \tilde{\Sigma}_\mu(x)=\left(\frac{\partial S_{\rm st}[\tilde{U}]}{
  \partial \tilde{U}_\mu(x)}\right)^T.
\end{equation}
This step depends on the form of the gauge and fermion action in terms
of the stout links and is usually just as straightforward as the
computation of $\Sigma_\mu^{\rm(th)}(x)$.  In subsequent steps,
$\Sigma^{(k)}$ and $U^{(k\!-\!1)}$ are used to compute
$\Sigma^{(k\!-\!1)}$, where the effective force at level $k$ is
defined by
\begin{equation}
  \Sigma_\mu^{(k)}(x)=\left(\frac{\partial S_{\rm st}[\tilde{U}]}{
  \partial U_\mu^{(k)}(x)}\right)^T.
\end{equation}
The recursive mapping
\begin{equation}
   \left\{  \Sigma^{(k)}, U^{(k-1)}\right\} 
  \longrightarrow \Sigma^{(k-1)}
 \label{eq:mapping}
\end{equation}
is repeated until $\Sigma^{(0)}$ is finally evaluated.

In Monte Carlo methods based on molecular dynamics, such as HMC\cite{hmc}
and the $R$-algorithm\cite{Ralgorithm}, the flow of lattice configurations
through phase space via Hamilton's equations is 
parametrized by a fictitious simulation time
coordinate $\tau$.  To determine the mapping in Eq.~(\ref{eq:mapping}),
it is convenient to express the force field in terms of the rate of change
of the action with respect to this time coordinate $\tau$:
\begin{equation}
  \frac{d}{d\tau}S_{\rm st}[\tilde{U}] = 2\sum_{x,\mu} \ReTr \left\{ 
        \Sigma_\mu^{(k)}(x)   \;
        \frac{d}{d\tau} U^{(k)}_\mu(x) 
\right\}, 
\end{equation}
for $k=0,1,\dots,n_\rho$. One then proceeds using the chain rule of
differentiation:
\begin{eqnarray}
\frac{dU^{(k)}_\mu(x)}{d\tau} &=&  e^{i Q^{(k-1)}_\mu(x)} 
   \frac{dU^{(k-1)}_\mu(x)}{d\tau} \nonumber\\
   &&+ \frac{d(e^{i Q^{(k-1)}_\mu(x)})}{d\tau} U^{(k-1)}_\mu(x).
\end{eqnarray}
To simplify matters, $Q^{(k-1)}_\mu(x),\ U^{(k-1)}_\mu(x),\ \Sigma^{(k-1)}_\mu(x)$
shall simply be written as $Q,U,\Sigma$, respectively,
in the calculations which follow, and $\Sigma^{(k)}_\mu(x)$ and
$U^{(k)}_\mu(x)$ shall be written as $\Sigma^\prime$ and $U^\prime$,
respectively.  The Cayley-Hamilton theorem
gives us
\begin{eqnarray}
    \frac{d(e^{i Q})}{d\tau} &=& \frac{d}{d\tau}
             \left(f_0 + f_1 Q + f_2 Q^2\right),\nonumber\\
   &=& \frac{df_0}{d\tau}+\frac{df_1}{d\tau}Q+\frac{df_2}{d\tau}Q^2\nonumber\\
    &&   + f_1\frac{dQ}{d\tau}+f_2\frac{dQ}{d\tau}Q+f_2Q\frac{dQ}{d\tau}.
\end{eqnarray}
Since the $f_j$ coefficients are functions $f_j=f_j(u,w)$ of $u$
and $w$ only, one has
\begin{equation}
    \frac{df_j}{d\tau} = \left(\frac{\partial f_j}{\partial u}\right)
   \frac{du}{d\tau}
    + \left(\frac{\partial f_j}{\partial w}\right)
   \frac{dw}{d\tau},
\end{equation}
and since $u$ and $w$ are functions of $c_0$ and $c_1$ only, then
\begin{eqnarray}
 \frac{du}{d\tau} &=& \left(\frac{\partial u}{\partial c_0}\right)
   \frac{dc_0}{d\tau}
                     +\left(\frac{\partial u}{\partial c_1}\right)
   \frac{dc_1}{d\tau},\\
 \frac{dw}{d\tau} &=& \left(\frac{\partial w}{\partial c_0}\right)
    \frac{dc_0}{d\tau}
                     +\left(\frac{\partial w}{\partial c_1}\right)
   \frac{dc_1}{d\tau}.
\end{eqnarray}
Next, one finds that
\begin{eqnarray}
    \frac{dc_0}{d\tau} &=& {\textstyle\frac{1}{3}}
    \frac{d}{d\tau} \Tr(Q^3) = 
 \Tr\left(Q^2 \frac{dQ}{d\tau}\right),\\
    \frac{dc_1}{d\tau} &=& {\textstyle\frac{1}{2}}
   \frac{d}{d\tau} \Tr (Q^2) = 
\Tr\left(Q \frac{dQ}{d\tau}\right),
\end{eqnarray}
and 
\begin{equation}
\begin{array}{ll}
\displaystyle\frac{\partial u}{\partial c_0}=\frac{1}{2(9u^2-w^2)}, &
\ \displaystyle \frac{\partial u}{\partial c_1}=\frac{u}{(9u^2-w^2)}, \\[5mm]
\displaystyle\frac{\partial w}{\partial c_0}=\frac{-3u}{2w(9u^2-w^2)}, &
\ \displaystyle \frac{\partial w}{\partial c_1}=\frac{3u^2-w^2}{2w(9u^2-w^2)}.
\end{array}
\end{equation}
Now define
\begin{equation}
r^{(1)}_j = \frac{\partial h_j}{\partial u},\qquad
r^{(2)}_j =\frac{1}{w}\frac{\partial h_j}{\partial w},
\end{equation}
and
\begin{eqnarray}
 b_{1j} &=& \frac{2ur_j^{(1)}+
    (3u^2\!-\!w^2)r^{(2)}_j-2(15u^2\!+\!w^2)f_j}{2(9u^2-w^2)^2},\\
 b_{2j}&=& \frac{r_j^{(1)}-3ur^{(2)}_j-24uf_j}{
   2(9u^2-w^2)^2},
\end{eqnarray}
then
\begin{equation}
 \frac{df_j}{d\tau}=b_{1j}\Tr\left(Q\frac{dQ}{d\tau}\right)
   +b_{2j}\Tr\left(Q^2\frac{dQ}{d\tau}\right),
\end{equation}
where
\begin{eqnarray}
r^{(1)}_0 &=& 2\Bigl(u+i(u^2\!-\!w^2)\Bigr)e^{2i u}\nonumber\\
  &&+2e^{-i u}\Bigl\{
 4u\bigl(2\!-\!i u  \bigr)\cos w \nonumber\\
  && +i\bigl(9u^2+w^2-i u( 3u^2+w^2) \bigr)\xi_0( w) 
  \Bigr\}, \\
r^{(1)}_1 &=& 2(1+2i u)e^{2i u}+e^{-i u}\Bigl\{
    -2(1-i u) \cos w \nonumber\\
  &&  +i\bigl( 6u +i(w^2-3u^2) \bigr)\xi_0( w) 
   \Bigr\}, \\
r^{(1)}_2 &=&  2i e^{2i u}+i e^{-i u}\Bigl\{
   \cos w -3 (1-i  u)  \xi_0( w)
   \Bigr\},\\
r^{(2)}_0 &=& -2e^{2i u}+2i ue^{-i u}\Bigl\{
  \cos w \nonumber\\
  && +( 1+4i u) \xi_0( w) 
  +3u^2\xi_1( w)\Bigr\}, \\
r^{(2)}_1 &=& -i e^{-i u}\Bigl\{
   \cos w  +(1+2i u)\xi_0( w)\nonumber\\ 
  &&-3u^2\xi_1( w)\Bigr\}, \\
r^{(2)}_2 &=&   e^{-i u}\Bigl\{
     \xi_0( w) 
   -3i  u \xi_1( w)\Bigr\},
\end{eqnarray}
with
\begin{eqnarray}
 \xi_0(w) &=& \frac{\sin w}{w},\\
 \xi_1(w) &=& \frac{\cos w}{w^2}-\frac{\sin w}{w^3}.
\end{eqnarray}
Given the above results,
\begin{eqnarray}
\frac{d(e^{i Q})}{d\tau}&=&\Tr\left(Q\frac{dQ}{d\tau}\right) B_1
+\Tr\left(Q^2\frac{dQ}{d\tau}\right) B_2\nonumber\\
&&+ f_1\frac{dQ}{d\tau}+f_2\frac{dQ}{d\tau}Q+f_2Q\frac{dQ}{d\tau},
\label{eq:dRdef}
\end{eqnarray}
where
\begin{equation}
 B_i = b_{i0}+b_{i1}\ Q + b_{i2}\ Q^2.
\end{equation}
The numerically sensitive $w\rightarrow 0$ limit has been
totally absorbed into $\xi_0$ and $\xi_1$, and
numerical problems as $c_0\rightarrow -c_0^{\rm max}$ can 
be circumvented by exploiting the following symmetry relation:
\begin{equation}
  b_{ij}(-c_0,c_1)=
  (-)^{i+j+1}\ b_{ij}^\ast(c_0,c_1). \label{eq:bsym}
\end{equation}
Once again, we emphasize that these coefficients are well-behaved,
non-singular functions of $c_0$ and $c_1$.

The rate of change of the stout links with respect to the simulation
time is then given by
\begin{eqnarray}
  \frac{dU^\prime}{d\tau} &=& e^{i Q} \frac{dU}{d\tau} 
  + \biggl\{  \Tr\left(Q\frac{dQ}{d\tau}\right) B_1
   +\Tr\left(Q^2\frac{dQ}{d\tau}\right) B_2 \nonumber\\
   &&\qquad + f_1\frac{dQ}{d\tau}+f_2\frac{dQ}{d\tau}Q+f_2Q\frac{dQ}{d\tau}
 \biggr\} U,
\end{eqnarray}
and hence, 
\begin{equation}
\ReTr\left(\Sigma^\prime
  \frac{dU^\prime}{d\tau}\right)
 =\ReTr\left(\Sigma^\prime e^{i Q} \frac{dU}{d\tau}\right)
                    - \ReTr\left(i\Lambda \frac{d\Omega}{d\tau}\right),
   \label{eqn:diff-stouttrace}
\end{equation}
defining 
\begin{eqnarray}
   \Lambda &=& 
   \frac{1}{2} (\Gamma + \Gamma^\dagger) - \frac{1}{2N} 
      \Tr\left(\Gamma + \Gamma^\dagger\right), \label{eq:Lambda}\\
  \Gamma &=& 
  \Tr(\Sigma^\prime B_1 U) \; Q         
+ \Tr(\Sigma^\prime B_2 U) \; Q^2 \nonumber\\
&&+ f_1  \; U \Sigma^\prime   
+ f_2  \; Q U \Sigma^\prime 
+ f_2  \; U \Sigma^\prime Q.
\label{eqn:gamma-def}
\end{eqnarray}
One sees that $\Gamma^{(k-1)}_\mu(x)$,
and hence, $\Lambda_\mu^{(k\!-\!1)}(x)$, are defined on each
lattice link in terms of the link variables $U^{(k-1)}_\mu(x)$ and 
$\Sigma^{(k)}_\mu(x)$. At this point, the details of the staple
construction in $C_\mu(x)$ must be included.  Given Eq.~(\ref{eq:Cdef})
and utilizing trace cyclicity and translational invariance,
one eventually obtains the following recursion relation:
\begin{eqnarray}
&&\Sigma_\mu(x) = \Sigma_\mu^\prime(x)
  \Bigl(f_0I+f_1Q_\mu(x)+f_2Q_\mu^2(x)\Bigr)\nonumber\\
   &&+iC^\dagger_\mu(x)\Lambda_\mu(x)  -i
 \displaystyle\sum_{\nu\neq\mu}\Bigl\{
\nonumber\\ &&
 \begin{array}[t]{l}
 \phantom{+} \rho_{\nu\mu}U_\nu(x\!+\!\hat{\mu}) 
       U^\dagger_\mu(x\!+\!\hat{\nu})U^\dagger_\nu(x)\Lambda_\nu(x) \\[1ex]
 +\rho_{\mu\nu}U^\dagger_\nu(x\!-\!\hat{\nu}\!+\!\hat{\mu})
   U^\dagger_\mu(x\!-\!\hat{\nu})\Lambda_\mu(x\!-\!\hat{\nu})
   U_\nu(x\!-\!\hat{\nu})    \\[1ex]
+\rho_{\nu\mu}U^\dagger_\nu(x\!-\!\hat{\nu}\!+\!\hat{\mu})
  \Lambda_\nu(x\!-\!\hat{\nu}\!+\!\hat{\mu})
  U^\dagger_\mu(x\!-\!\hat{\nu}) U_\nu(x\!-\!\hat{\nu}) \\[1ex]
-\rho_{\nu\mu} U^\dagger_\nu(x\!-\!\hat{\nu}\!+\!\hat{\mu})
 U^\dagger_\mu(x\!-\!\hat{\nu}) \Lambda_\nu(x\!-\!\hat{\nu})
  U_\nu(x\!-\!\hat{\nu}) \\[1ex]
-\rho_{\nu\mu}  \Lambda_\nu(x\!+\!\hat{\mu})U_\nu(x\!+\!\hat{\mu})
   U^\dagger_\mu(x\!+\!\hat{\nu})
   U^\dagger_\nu(x) \\[1ex]
+\rho_{\mu\nu}U_\nu(x\!+\!\hat{\mu})U^\dagger_\mu(x\!+\!\hat{\nu})
 \Lambda_\mu(x\!+\!\hat{\nu})
 U^\dagger_\nu(x) 
 \Bigr\},\end{array} \nonumber\\[-6mm]
\label{eq:sigma}
\end{eqnarray}
in which unprimed quantities refer to step $(k-1)$ and primed quantities
refer to step $k$.

Although the force field is related to the underlying link variables in
a complicated manner, each step in the recursive computation
of the force field outlined above is straightforward, facilitating a
natural and efficient implementation in software.  Analyticity in the
entire (finite) complex plane is an important property here; if 
$\Sigma^{(k)}_\mu(x)$ is well behaved, then $\Sigma^{(0)}_\mu(x)$ must
be also.  Note that the above formalism can be easily adapted for other
definitions of $C_\mu(x)$.

